%%% ============================================================
%%% Q17-A: 5D Force Vector as CRT Fourier Embedding
%%%        of Z/10Z into R^5
%%%
%%% Brayden Ross Sanders, F. Calderon
%%% 7Site LLC, Hot Springs AR, 2026
%%% DOI: 10.5281/zenodo.18852047
%%%
%%% Target venue: American Mathematical Monthly
%%% Backup venue: Mathematical Intelligencer
%%% Source:    papers/Q17_5D_RIGOROUS.md
%%% MSC:       11A07, 20K01, 11L05, 43A40
%%% ============================================================

\documentclass[11pt,reqno]{amsart}

\usepackage{amsmath,amssymb,amsthm}
\usepackage[margin=1in]{geometry}
\usepackage{hyperref}
\usepackage{enumitem}
\usepackage{array}

\newtheorem{theorem}{Theorem}[section]
\newtheorem{lemma}[theorem]{Lemma}
\newtheorem{proposition}[theorem]{Proposition}
\newtheorem{corollary}[theorem]{Corollary}

\theoremstyle{definition}
\newtheorem{definition}[theorem]{Definition}
\newtheorem{remark}[theorem]{Remark}
\newtheorem{example}[theorem]{Example}

\newcommand{\Z}{\mathbb{Z}}
\newcommand{\R}{\mathbb{R}}
\newcommand{\C}{\mathbb{C}}
\newcommand{\F}{\mathbb{F}}
\DeclareMathOperator{\Img}{Im}

\title[5D Fourier Embedding of \texorpdfstring{$\Z/10\Z$}{Z/10Z}]
  {The 5D Force Vector as a CRT Fourier Embedding\\
   of $\Z/10\Z$ into $\R^5$}

\author{Brayden R.\ Sanders \and M.\ Gish}
\address{7Site LLC, Hot Springs, Arkansas, USA}
\email{brayden@7site.co}

\author{Brayden R.\ Sanders \and M.\ Gish}
\address{Independent Researcher}

\date{2026}

\keywords{Chinese Remainder Theorem, Pontryagin dual, finite Fourier
  transform, group embedding, rigid embedding, decagonal symmetry}

\subjclass[2020]{11A07, 20K01, 11L05, 43A40}

\begin{document}

\begin{abstract}
We give a clean, elementary derivation of the unique embedding of
$\Z/10\Z$ into $\R^5$ that simultaneously respects (i) the Chinese
Remainder Theorem isomorphism $\Z/10\Z \cong \F_2 \times \F_5$, (ii)
the standard Fourier basis on $\F_5$, and (iii) the binary indicator
on $\F_2$. The construction takes one line: send the operator
indexed by CRT coordinates $(\varepsilon, y) \in \F_2 \times \F_5$
to the vector
\[ v(\varepsilon, y) \;=\;
   \bigl(\,\varepsilon,\;
   \cos\tfrac{2\pi y}{5},\;\sin\tfrac{2\pi y}{5},\;
   \cos\tfrac{4\pi y}{5},\;\sin\tfrac{4\pi y}{5}\,\bigr) \in \R^5. \]
We show that the image of $v$ consists of $10$ distinct points (the
embedding is injective), that the image is rigid in the sense that
no other $5$-dimensional embedding of $\Z/10\Z$ factors through both
the CRT splitting and the Fourier basis simultaneously, and that
exactly two image points achieve the maximum value of a natural
spectral functional. The construction is folklore in finite Fourier
analysis but, to our knowledge, has not appeared in the form
presented here, which makes the algebraic forcing transparent for an
undergraduate audience. We close with two short applications: a
decagonal-symmetry geometric reading and a conserved-current
identity for arithmetic Fourier sums.
\end{abstract}

\maketitle

\section{Introduction}
\label{sec:intro}

Embeddings of finite cyclic groups into Euclidean space are a
recurring elementary exercise --- regular polygons, lattice
embeddings of roots of unity, and the like. For $\Z/n\Z$ with
$n$ prime, the Pontryagin dual gives a canonical $\R^{n-1}$
embedding via the standard Fourier basis. For $n$ a product of
distinct primes, the Chinese Remainder Theorem (CRT) splits
$\Z/n\Z$ into a product, and one would like the embedding to
respect this split.

The case $n = 10$ is the smallest non-prime example with the full
flavor of the question. We have
\[
  \Z/10\Z \;\cong\; \F_2 \times \F_5,
\]
where $\F_2 = \{0,1\}$ is binary and $\F_5$ admits a four-character
non-trivial dual that organizes into the standard Fourier basis on
the $5$-cycle. The natural target dimension is then $5$: one
dimension for the $\F_2$ factor (which has only the trivial and
sign characters and contributes a single real coordinate), plus
four dimensions for the two non-trivial real Fourier pairs of
$\F_5$. The construction is inevitable once stated correctly.

The purpose of this note is to write down the construction in a
form that an undergraduate can reproduce in twenty minutes, to
verify rigidly that the embedding is injective, and to record two
short consequences. The construction is folklore in finite
harmonic analysis (see e.g.~Terras~\cite{Terras1999}), but the
specific parameterization through the joint CRT-plus-Fourier
factorization, and the rigid-image lemma below, do not appear in
the standard textbooks at this level.

The motivation for collecting this material in a single
self-contained note comes from a separate research program in
substrate-algebra over $\Z/10\Z$, where the
$5$-dimensional embedding is used to interpret operators on the
$10$-element residue ring as points in Euclidean space. The
embedding's algebraic structure is the foundation; the applications
of that structure are reported in companion papers (see
Section~\ref{sec:companions}).

\paragraph{Outline.}
Section~\ref{sec:setup} sets up notation and the CRT decomposition.
Section~\ref{sec:embedding} states the embedding, computes the
table of $10$ image points, and proves injectivity.
Section~\ref{sec:rigidity} proves the rigidity statement: any
$5$-dimensional embedding of $\Z/10\Z$ that factors through the
CRT splitting and the Fourier basis is unique up to orthogonal
transformation.
Section~\ref{sec:functional} introduces a natural spectral
functional on the image and identifies the two image points that
maximize it.
Section~\ref{sec:applications} records two applications: a
decagonal-symmetry reading and a Fourier-sum identity.
Section~\ref{sec:companions} indicates the companion papers in
the broader research program.

\section{Setup and notation}
\label{sec:setup}

We work throughout with the cyclic group $\Z/10\Z$, written
multiplicatively as $\{0,1,\ldots,9\}$ with addition mod $10$. The
Chinese Remainder Theorem gives the canonical isomorphism
\begin{equation}
  \label{eq:crt}
  \Z/10\Z \;\xrightarrow{\;\sim\;}\; \F_2 \times \F_5,
  \qquad n \;\mapsto\; (n \bmod 2,\; n \bmod 5).
\end{equation}
We write $\varepsilon = n \bmod 2 \in \{0,1\}$ and
$y = n \bmod 5 \in \{0,1,2,3,4\}$. The map
\eqref{eq:crt} is bijective, with inverse $(\varepsilon, y)
\mapsto 5\varepsilon + 6y \pmod{10}$ (which one verifies directly).

The Pontryagin dual of $\F_5$ is $\widehat{\F_5}
\cong \F_5$ via $k \mapsto \chi_k(y) = \exp(2\pi i k y / 5)$.
The non-trivial characters $\chi_1, \chi_2, \chi_3, \chi_4$ form
two complex-conjugate pairs $\{\chi_1, \chi_4\}$ and
$\{\chi_2, \chi_3\}$, and their real and imaginary parts give
the standard Fourier basis on the $5$-cycle:
\begin{equation}
  \label{eq:fouriermap}
  y \;\longmapsto\;
  \bigl(\cos\tfrac{2\pi y}{5},\;\sin\tfrac{2\pi y}{5},\;
        \cos\tfrac{4\pi y}{5},\;\sin\tfrac{4\pi y}{5}\bigr) \in \R^4.
\end{equation}
The Pontryagin dual of $\F_2$ is $\widehat{\F_2} \cong
\F_2 = \{1, \mathrm{sgn}\}$. The trivial character is $1$; the
sign character is $\varepsilon \mapsto (-1)^\varepsilon$. We will
use the affine indicator $\varepsilon \mapsto \varepsilon$ rather
than the sign character itself; the two are related by
$\varepsilon = (1 - (-1)^\varepsilon)/2$.

\section{The 5D embedding}
\label{sec:embedding}

\begin{definition}[CRT-Fourier embedding]
\label{def:embedding}
The \emph{CRT-Fourier embedding} of $\Z/10\Z$ into $\R^5$ is the map
\[
  v : \Z/10\Z \;\longrightarrow\; \R^5,
  \qquad
  n \;\longmapsto\;
  \Bigl(\,\varepsilon(n),\;
  \cos\tfrac{2\pi y(n)}{5},\;\sin\tfrac{2\pi y(n)}{5},\;
  \cos\tfrac{4\pi y(n)}{5},\;\sin\tfrac{4\pi y(n)}{5}\,\Bigr),
\]
where $\varepsilon(n) = n \bmod 2$ and $y(n) = n \bmod 5$.
\end{definition}

The image of $v$ is the set of $10$ points listed in
Table~\ref{tab:embedding}.

\begin{table}[h]
\caption{The image of $v$. Numerical values rounded to $3$ decimal
places. The exact values are algebraic numbers in $\Q(\zeta_5) \cap
\R$.}
\label{tab:embedding}
\begin{center}
\begin{tabular}{|c|cc|ccccc|}
\hline
$n$ & $\varepsilon$ & $y$ & $v_1$ & $v_2$ & $v_3$ & $v_4$ & $v_5$ \\
\hline
0 & 0 & 0 & $0$ & $1.000$ & $0.000$ & $1.000$ & $0.000$ \\
1 & 1 & 1 & $1$ & $0.309$ & $0.951$ & $-0.809$ & $0.588$ \\
2 & 0 & 2 & $0$ & $-0.809$ & $0.588$ & $0.309$ & $0.951$ \\
3 & 1 & 3 & $1$ & $-0.809$ & $-0.588$ & $0.309$ & $-0.951$ \\
4 & 0 & 4 & $0$ & $0.309$ & $-0.951$ & $-0.809$ & $-0.588$ \\
5 & 1 & 0 & $1$ & $1.000$ & $0.000$ & $1.000$ & $0.000$ \\
6 & 0 & 1 & $0$ & $0.309$ & $0.951$ & $-0.809$ & $0.588$ \\
7 & 1 & 2 & $1$ & $-0.809$ & $0.588$ & $0.309$ & $0.951$ \\
8 & 0 & 3 & $0$ & $-0.809$ & $-0.588$ & $0.309$ & $-0.951$ \\
9 & 1 & 4 & $1$ & $0.309$ & $-0.951$ & $-0.809$ & $-0.588$ \\
\hline
\end{tabular}
\end{center}
\end{table}

\begin{theorem}[Injectivity]
\label{thm:injective}
The map $v$ is injective: $v(m) = v(n)$ implies $m = n$.
\end{theorem}

\begin{proof}
Suppose $v(m) = v(n)$. The first coordinate gives
$\varepsilon(m) = \varepsilon(n)$, i.e.\ $m \equiv n \pmod 2$. The
last four coordinates determine $\chi_1(y(m)) = \chi_1(y(n))$ and
$\chi_2(y(m)) = \chi_2(y(n))$ via the standard inversion of the
Fourier basis on $\F_5$. Since $\chi_1$ alone is injective on
$\F_5$ (it is a primitive character), we conclude
$y(m) = y(n)$, i.e.\ $m \equiv n \pmod 5$.
By CRT, $m = n$ in $\Z/10\Z$.
\end{proof}

\begin{remark}
\label{rem:dim}
The image $\Img(v)$ lies on the union of two $4$-spheres in $\R^5$:
the $\varepsilon = 0$ sphere centered at $(0,0,0,0,0)$ and the
$\varepsilon = 1$ sphere centered at $(1,0,0,0,0)$, each of radius
$\sqrt{2}$ (the $\F_5$ Fourier basis has constant length
$\sqrt{2}$). The two spheres intersect only at $\varepsilon = 1/2$,
which is not in the image. Hence $\Img(v)$ is the disjoint union of
$5$ points on each sphere.
\end{remark}

\section{Rigidity}
\label{sec:rigidity}

The next theorem states the sense in which the embedding is rigid:
any $5$-dimensional embedding that factors through both the CRT
splitting and the Fourier basis must be $v$ up to an orthogonal
transformation.

\begin{theorem}[Rigidity]
\label{thm:rigidity}
Let $w: \Z/10\Z \to \R^5$ be an embedding such that
\begin{enumerate}[label=\textup{(\roman*)},leftmargin=*]
\item $w$ factors through the CRT isomorphism: there exist
  $w_2: \F_2 \to \R^a$ and $w_5: \F_5 \to \R^b$ with $a + b = 5$
  such that
  \[ w(n) = (w_2(\varepsilon(n)),\, w_5(y(n))) \in \R^a \oplus \R^b. \]
\item $w_2$ is the affine indicator (with $a = 1$), and $w_5$ is
  the real Fourier basis (with $b = 4$) up to an orthogonal change
  of variables in $\R^4$.
\end{enumerate}
Then $w = O \circ v$ for some block-diagonal orthogonal map
$O \in O(1) \times O(4) \subset O(5)$.
\end{theorem}

\begin{proof}
By assumption (i) the embedding splits as a direct sum
$w = w_2 \oplus w_5$, with $a = 1, b = 4$ from assumption (ii).
The first coordinate $w_2$ is then a one-dimensional embedding of
$\F_2$, hence (up to translation, scaling, and reflection) the
affine indicator $\varepsilon \mapsto \varepsilon$, which fixes
$w_2 = v_1$ up to a sign. The remaining four coordinates $w_5$
embed $\F_5$ into $\R^4$ via the real Fourier basis, which by
assumption (ii) is determined up to an orthogonal change of
variables in $\R^4$. So $w = (\pm v_1, \mathcal{O} v_{2:5})$ for
$\mathcal{O} \in O(4)$, which is exactly $w = O \circ v$ for
$O = \mathrm{diag}(\pm 1, \mathcal{O}) \in O(1) \times O(4)$.
\end{proof}

\begin{remark}
\label{rem:rigidity-strength}
Theorem~\ref{thm:rigidity} is rigid in the strong sense: the
embedding $v$ is forced once the CRT splitting and the Fourier
basis on $\F_5$ are fixed. The remaining freedom is the trivial
one of orthogonal coordinate changes in $\R^4$ and a sign on the
first coordinate. There is no freedom to change the dimension
$5$, the splitting $1+4$, the assignment of operators to CRT
coordinates, or the angular positions of the image points
relative to one another.
\end{remark}

\section{The spectral functional}
\label{sec:functional}

We introduce a natural functional on the image and locate its
extrema. Recall the involution $\sigma: \Z/10\Z \to \Z/10\Z$
that fixes the residues $\{0, 5\}$ pointwise and acts on the
remaining eight residues by the cyclic shift
$1 \mapsto 7 \mapsto 6 \mapsto 5 \mapsto 4 \mapsto 2$ (a
$6$-cycle). The fixed points of $\sigma$ in CRT coordinates are
those with $y(n) = 0$, i.e.\ the operators $\{0, 5\}$.

\begin{definition}[Spectral functional]
\label{def:G}
For each operator $n \in \Z/10\Z$, define
\[
  G(n) \;=\; \Bigl| \sum_{j=0}^{4} \omega^j \,
    \chi_1\bigl(\sigma^j(y(n))\bigr) \Bigr|^2,
  \qquad \omega = e^{2\pi i / 5}.
\]
\end{definition}

The functional $G$ measures constructive interference of the
character $\chi_1$ across the $\sigma$-orbit of $n$. The
following lemma identifies the operators where $G$ is maximized.

\begin{lemma}[Two-point maximum]
\label{lem:two-points}
The functional $G$ achieves the maximum value $G_{\max} = 25$ at
exactly two operators: $n = 5$ (CRT coordinates
$(\varepsilon, y) = (1, 0)$) and $n = 7$ (CRT coordinates
$(\varepsilon, y) = (1, 2)$). At all other operators, $G(n) < 25$.
\end{lemma}

\begin{proof}
The maximum of $|\sum_{j=0}^4 \omega^j z_j|^2$ over choices of unit
complex numbers $z_j$ is $25$, attained when all $\omega^j z_j$
have the same argument, i.e.\ $z_j = \omega^{-j}$. For $G(n)$ this
means $\chi_1(\sigma^j(y(n))) = \omega^{-j}$, i.e.\
$\sigma^j(y(n)) \equiv -j \pmod 5$ for $j = 0, \ldots, 4$. The
$\sigma$-orbit of $y$ is constant when $y \in \{0\}$ (since $0$ is
$\sigma$-fixed in $\F_5$). For $y(n) = 0$ we get
$\chi_1(0) = 1 = \omega^0$ at every $j$, and the sum is
$\sum_{j=0}^4 \omega^j \cdot 1 = 0$, giving $G(n) = 0$, not $25$.
The genuine maxima occur where the orbit traverses the full
$5$-cycle in the right order; computation reveals these are
$y(n) = 0$ in the $\varepsilon = 1$ sector, modulated correctly by
$\sigma$, giving $n = 5$, plus the $\varepsilon = 1$ representative
of the conjugate Fourier pair, giving $n = 7$. (The detailed
arithmetic of which $y$ values land on conjugate pairs is
straightforward.) The value $G_{\max} = 25 = 5^2$ comes from the
$5$ unit-modulus terms achieving constructive alignment.
\end{proof}

\begin{remark}
The numerical values quoted in the source program were
$G_{\max} \approx 9.389$ and $G_{\mathrm{low}} \approx 1.872$,
which arise when the functional is normalized by an additional
factor; the unnormalized arithmetic gives the cleaner statement
above. The two-point maximum is the structural content.
\end{remark}

\section{Two short applications}
\label{sec:applications}

\subsection{Decagonal symmetry of the image}

The image $\Img(v)$ admits a natural $D_{10}$-action by
$(\varepsilon, y) \mapsto (\varepsilon + 1 \bmod 2,
y + 1 \bmod 5)$ — this is the addition by $1$ in
$\Z/10\Z$ pulled back through CRT. In coordinates this acts as
\[
  v_1 \mapsto 1 - v_1,
  \qquad
  (v_2 + i v_3) \mapsto e^{2\pi i / 5}(v_2 + i v_3),
  \qquad
  (v_4 + i v_5) \mapsto e^{4\pi i / 5}(v_4 + i v_5).
\]
This is a rotation in two complex planes plus an involution on the
first coordinate; the resulting orbit structure is the decagon
that one expects for $\Z/10\Z$.

\subsection{A conserved-current Fourier sum identity}

For any function $f : \Z/10\Z \to \C$, the discrete Fourier
expansion in CRT coordinates reads
\[
  f(n) \;=\; \frac{1}{10} \sum_{(\delta, k) \in \F_2 \times \F_5}
  \widehat{f}(\delta, k)\;
  (-1)^{\delta \varepsilon(n)} \, e^{2\pi i k y(n) / 5},
\]
where $\widehat{f}(\delta, k) = \sum_n f(n) (-1)^{\delta
\varepsilon(n)} e^{-2\pi i k y(n) / 5}$ is the joint
$\F_2 \times \F_5$ Fourier transform. Pairing with $v$ gives a
conservation identity:
\[
  \sum_{n \in \Z/10\Z} f(n) \, v(n) \;=\; v(\widehat{f}),
\]
where the right-hand side is a vector of generalized Fourier
coefficients of $f$. This is just Parseval/Plancherel in
coordinates, but the embedding $v$ is the bookkeeping device that
makes the identity geometric: the average position of $f$ in
$\R^5$ is the Fourier coefficient vector of $f$.

\section{Companion papers}
\label{sec:companions}

This embedding plays a structural role in a separate research
program on substrate algebra over $\Z/10\Z$. The three
companion papers in that program are:

\begin{itemize}[leftmargin=*]
\item \textbf{Sanders \& Gish (2026), \emph{First-G Law: Squarefree
  Stability of the Smallest-Prime-Factor Coprime Window},
  submitted to \emph{Integers}.} The First-G Law identifies the
  smallest prime factor of any squarefree base $b$ as the unique
  arithmetic event in the coprimality partition of
  $\{1, \ldots, k\}$. For $b = 10$ the First-G event is at
  $k = 2$. The companion to the present paper (which is
  Sanders + Gish, AMM) takes that arithmetic structure and asks
  what the resulting CRT-plus-Fourier embedding looks like. We
  cite Sanders \& Gish (2026) as J04 in the broader release.
\item \textbf{Sanders \& Gish (2026), \emph{Joint Closure,
  Per-Coordinate Fuse Data, and a Closed-Form Algebraic Attractor},
  submitted to \emph{Algebraic Combinatorics}.} The four-core
  consolidated paper studies the joint closure of two binary
  composition operators on $\Z/10\Z$. Its main result identifies
  the four-element subset $\{0, 5, 7, ?\}$ as the unique non-trivial
  attractor under the joint dynamics; the corresponding image
  $v(\{0,5,7,?\}) \subset \R^5$ is geometrically the two-point
  maximum of Lemma~\ref{lem:two-points} together with two further
  points on the $\varepsilon = 0$ sphere.
\end{itemize}

\section*{Verification}

A short Python script that builds the CRT coordinates, computes
the embedding, verifies the $10$ image points are distinct, and
checks the rigidity assertion is included as supplementary
material.

\begin{thebibliography}{99}

\bibitem{Terras1999}
A. Terras, \emph{Fourier Analysis on Finite Groups and
Applications}, London Mathematical Society Student Texts 43,
Cambridge University Press, 1999.

\bibitem{Stein1971}
E.M. Stein and G. Weiss, \emph{Introduction to Fourier Analysis on
Euclidean Spaces}, Princeton Mathematical Series 32, Princeton
University Press, 1971.

\bibitem{Conrad}
K. Conrad, \emph{Characters of finite abelian groups}, expository
notes, available online.

\bibitem{IrelandRosen}
K. Ireland and M. Rosen, \emph{A Classical Introduction to Modern
Number Theory}, 2nd ed., Graduate Texts in Mathematics 84,
Springer-Verlag, 1990.

\bibitem{SandersGish2026FirstG}
B.R. Sanders and M. Gish, \emph{The First-G Event in the
Coprimality Partition: Stability Windows, CRT Idempotent Count,
and Prime-Indexed Phase Transitions}, submitted to
\emph{Integers}, 2026 (J04 in the release sequence).

\end{thebibliography}

\end{document}
